\BourbakiExercises{Exercises}

\begin{enumerate}
\def\labelenumi{\arabic{enumi})}
\tightlist
\item
  \begin{enumerate}
  \def\labelenumii{\alph{enumii})}
  \tightlist
  \item
    Let M be an A-module and x an element of M whose annihilator is the intersection of a finite family of maximal ideals. Prove that the module Ax is semisimple (embed it in a finite product of simple modules).
  \end{enumerate}
\end{enumerate}

\begin{enumerate}
\def\labelenumi{\alph{enumi})}
\setcounter{enumi}{1}
\tightlist
\item
  Deduce that a module M is semisimple if and only if the annihilator of every nonzero element is the intersection of a finite family of maximal ideals.
\end{enumerate}

\begin{enumerate}
\def\labelenumi{\arabic{enumi})}
\setcounter{enumi}{1}
\tightlist
\item
  Let M be a module. Prove that the following properties are equivalent:
\end{enumerate}

\(\alpha)\) M is semisimple and has finite length.

\(\beta)\) M is Noetherian, and every maximal submodule of M admits a supplement.

\(\gamma)\) M is Artinian, and every simple submodule of M admits a supplement.

\begin{enumerate}
\def\labelenumi{\arabic{enumi})}
\setcounter{enumi}{2}
\item
  Let \(\mathrm{K}\) be a field, I an infinite set, and A the ring \(\mathrm{K}^{\mathrm{I}}\) . Prove that the A-module \(\mathrm{A}_s\) is not semisimple even though every monogenous ideal admits a supplement (determine all simple submodules of \(\mathrm{A}_s\) ).
\item
  Let M be a module that is the direct sum of a family \((\mathrm{M}_{\alpha})_{\alpha\in\mathrm{I}}\) of submodules. Prove that the submodules \(M_{\alpha}\) are stable under all automorphisms of M if and only if every homomorphism \(M_{\alpha}\to M_{\beta}\) with \(\alpha\neq\beta\) is zero.
\item
  Let M be a module of finite length, and let \(M = \oplus_{i \in I} M_i\) be a decomposition of M into a direct sum of indecomposable submodules. Let \(J \subset I\) be an equivalence class for the relation ``the modules \(M_i\) and \(M_j\) are isomorphic'' between i and j; we set \(M_J = \oplus_{i \in J} M_i\) . Give an example to show that \(M_J\) is not necessarily stable under every automorphism of M (use Exercise 4 with M taken to be a Z-module of finite length).
\item
  A module M is called semi-Artinian if every nonzero quotient of M contains a simple submodule. Every Artinian module and every vector space over a field is semi-Artinian.
\end{enumerate}

\begin{enumerate}
\def\labelenumi{\alph{enumi})}
\item
  Let \(0 \rightarrow M' \rightarrow M \rightarrow M'' \rightarrow 0\) be an exact sequence. Then M is semi-Artinian if and only if the same holds for \(M'\) and \(M''\) (if M is semi-Artinian, consider a maximal element in the set of submodules N of M such that \(N \cap M' = 0\) ).
\item
  The direct sum of a family of semi-Artinian modules is a semi-Artinian module. c) Let K be a field and A the product ring \(K^{N}\) . Prove that the A-module \(A_{s}\) is not semi-Artinian even though it is a product of simple A-modules (if s is the socle of \(A_{s}\) , prove that the A-module \(A_{s}/s\) does not contain any simple modules).
\item
  Let A be a left Noetherian ring and M a semi-Artinian A-module. Prove that M is the sum of its Artinian submodules (let S be the sum of the Artinian submodules of M; observe that if N is a finitely generated submodule of M with simple image in M/S, then \(N \cap S\) is Artinian). In particular, if M is finitely generated, then it is Artinian.
\item
  Let D be a field, V an infinite-dimensional D-vector space, and A the subring of \(\operatorname{End}_{\mathrm{D}}(\mathrm{V})\) generated by the center of D and the endomorphisms of finite rank. Prove that \(A_{s}\) is semi-Artinian (use a) and b)) but is not the sum of its Artinian submodules.
\end{enumerate}

\begin{enumerate}
\def\labelenumi{\arabic{enumi})}
\setcounter{enumi}{6}
\tightlist
\item
  Let A be a ring and M an A-module. A submodule N of M is called essential if \(N \cap P \neq 0\) for every nonzero submodule P of M.
\end{enumerate}

\begin{enumerate}
\def\labelenumi{\alph{enumi})}
\item
  Prove that every submodule N of M is a direct factor of an essential submodule (consider a submodule P that is maximal among the submodules of A with zero intersection with N).
\item
  Prove that the socle of M is the intersection of the essential submodules of M (deduce from a) that this intersection is semisimple). In particular, M is semisimple if and only if it does not contain any essential submodules other than itself.
\end{enumerate}

¶ 8) Let A be a left Noetherian ring.

\begin{enumerate}
\def\labelenumi{\alph{enumi})}
\item
  Let x be a nonzero element of A. Prove that the left annihilator Ann x of x is not essential (among the elements that do not have this property, take an x such that Ann x is maximal, and consider \(Ax \cap Ann x\) ).
\item
  Let x be a right cancellable element of A. Prove that the ideal Ax is essential (if a is an ideal of A such that \(a \cap Ax = 0\) , consider the ideal \(a + ax + ax^{2} + \ldots\) ).
\item
  Deduce from a) and b) that every right cancellable element of A is cancellable.
\item
  Let a be an essential left ideal of A containing a nonnilpotent element. Prove that a contains a cancellable element (let \((x_{1},\ldots,x_{n})\) be a sequence of nonzero elements of a satisfying Ann \(x_{i}^{2}=\) Ann \(x_{i}\) and \(x_{i+1}\in\) Ann \(x_{i}\cap\cdots\cap\) Ann \(x_{1}\) for every \(i\geqslant1\) ; observe that a contains the direct sum of the \(Ax_{i}\) . Deduce that there exists such a sequence with \(\bigcap_{i}\) Ann \(x_{i}=0\) ; prove that the element \(\sum x_{i}^{2}\) is then right cancellable, and use c)).
\end{enumerate}

